\section{Dynamic Optimal Transport}
\label{sec-dynamic-optimal-transport}
\subsection{Evolutions over the Space of Measures}
\paragraph{Lagrangian and Eulerian descriptions.}
\begin{equation}
\frac{\d x(t)}{\d t} = v_t(x(t)), \label{eq:lagrangian-advection}
\end{equation}
\begin{equation}
\frac{\partial \alpha_t}{\partial t} + \diverg(v_t \alpha_t) = 0. \label{eq:eulerian-advection}
\end{equation}
\begin{equation}\label{eq:eulerian-advection-weak}
\int_0^1\!\int_{\RR^d}
\left(\partial_t\varphi(t,x)+\dotp{v_t(x)}{\nabla_x\varphi(t,x)}\right)
\d\alpha_t(x)\d t
=
0.
\end{equation}
This equation is obtained from~\eqref{eq:eulerian-advection} by integration by parts. Hence, for smooth positive densities, the classical and weak formulations are equivalent; the weak formulation is useful because it still makes sense for discrete measures whose particles evolve according to~\eqref{eq:lagrangian-advection}.

\begin{prop}[Lagrangian flows solve the continuity equation]\label{prop-lagrangian-flow-continuity}
Let $(T_t)_{t\in[0,1]}$ be a $C^1$ family of diffeomorphisms of $\RR^d$ and define $\alpha_t=(T_t)_\sharp\alpha_0$. Assume that the derivatives below are integrable with respect to $\alpha_0$, and define the Eulerian velocity field by
\(v_t(T_t(y))=\partial_t T_t(y).\)
Then $(\alpha_t,v_t)$ solves the continuity equation in the weak sense~\eqref{eq:eulerian-advection-weak}.
In particular, if $\alpha_0=\frac1n\sum_i\delta_{x_i(0)}$ is empirical, then $\alpha_t=\frac1n\sum_i\delta_{x_i(t)}$ is empirical as well, with particle velocities $\dot x_i(t)=v_t(x_i(t))$.
\end{prop}
\begin{proof}
Let $\varphi\in C_c^1((0,1)\times\RR^d)$. Since $\alpha_t=(T_t)_\sharp\alpha_0$, \(\frac{\d}{\d t}\int \varphi(t,x)\d\alpha_t(x) = \frac{\d}{\d t}\int \varphi(t,T_t(y))\d\alpha_0(y).\)
The chain rule gives
\[
\frac{\d}{\d t}\int \varphi(t,T_t(y))\d\alpha_0(y)
=
\int
\left(
\partial_t\varphi(t,T_t(y))
+\dotp{\nabla_x\varphi(t,T_t(y))}{\partial_t T_t(y)}
\right)
\d\alpha_0(y).
\]
Using the definition of $v_t$ and the push-forward relation, this equals
\[
\int
\left(
\partial_t\varphi(t,x)+\dotp{\nabla_x\varphi(t,x)}{v_t(x)}
\right)
\d\alpha_t(x).
\]
Integrating in time and using the boundary values of $\varphi$ gives~\eqref{eq:eulerian-advection-weak}.
\end{proof}
\paragraph{From measure evolutions to vector fields.}
For a given evolution $(\alpha_t)_t$, there are typically infinitely many velocity fields $v_t$ satisfying

\begin{equation}\label{eq:inverse-flow}
\partial_t\alpha_t+\diverg(\alpha_t v_t)=0.
\end{equation}
\(\Hh_\alpha=\enscond{v\in L^2(\alpha;\RR^d)}{\diverg(\alpha v)=0\ \text{in distributions}}.\)
\paragraph{Dacorogna--Moser inversion.}
\begin{equation}\label{eq:dacorogna-moser}
v_t=-\frac{1}{\rho_t}\nabla\Delta^{-1}(\partial_t\rho_t),
\end{equation}
\(\alpha_t=(1-t)\alpha_0+t\alpha_1=\rho_t\,\d x, \qquad \rho_t=(1-t)\rho_0+t\rho_1.\)
\(\diverg w=\rho_0-\rho_1, \qquad w\cdot n=0\quad\hbox{on }\partial\Omega,\)
for instance $w=-\nabla\phi$ with $\Delta\phi=\rho_1-\rho_0$ and Neumann boundary condition. Then

\(v_t=\frac{w}{\rho_t}\)
satisfies $\partial_t\rho_t+\diverg(\rho_t v_t)=0$. The flow $\partial_t T_t=v_t\circ T_t$, $T_0=\Id$, therefore transports $\rho_0\d x$ onto $\rho_t\d x$, and $T_1$ solves the prescribed-Jacobian problem $\rho_1(T_1(x))\det(\nabla T_1(x))=\rho_0(x)$.

\paragraph{Least-square inversion and gradient structure.}
\begin{equation}\label{eq:least-square-field}
\umin{v}
\frac12\int_0^1\!\int_{\RR^d}\norm{v_t(x)}^2\,\d\alpha_t(x)\d t
\quad
\text{subject to}
\quad
\partial_t\alpha_t+\diverg(\alpha_t v_t)=0.
\end{equation}
\begin{prop}[Least-square velocities are gradients]\label{prop-least-square-gradient-velocity}
Assume that $\alpha_t=\rho_t\,\d x$ is a smooth positive density curve, that $\partial_t\rho_t$ has zero integral, and that boundary terms vanish. The minimizer of~\eqref{eq:least-square-field}, if it exists, is a gradient field
\(v_t=\nabla\phi_t,\)
where $\phi_t$, unique up to an additive constant on each connected component, solves the weighted Poisson equation
\begin{equation}\label{eq:least-square-field-explicit}
-\diverg(\rho_t\nabla\phi_t)=\partial_t\rho_t,
\qquad
v_t=-\nabla\Delta_{\alpha_t}^{-1}(\partial_t\alpha_t),
\quad
\Delta_{\alpha_t}\phi=\diverg(\alpha_t\nabla\phi).
\end{equation}
\end{prop}
\begin{proof}
Introduce a Lagrange multiplier $\phi_t$ for the continuity equation. The constrained problem has the formal saddle formulation
\[
\umin{v}\umax{\phi}
\int_0^1
\left[
\frac12\int_{\RR^d}\norm{v_t(x)}^2\,\d\alpha_t(x)
+\int_{\RR^d}\phi_t(x)\big(\diverg(\alpha_t v_t)(x)+\partial_t\alpha_t(x)\big)\d x
\right]\d t.
\]
Integrating by parts in the divergence term gives, for each $t$,
\[
\int \left(\frac12\norm{v_t}^2-\dotp{\nabla\phi_t}{v_t}\right)\d\alpha_t
+
\int \phi_t\,\partial_t\alpha_t.
\]
The pointwise minimizer in $v_t$ is therefore $v_t=\nabla\phi_t$. Substituting this into the constraint $\partial_t\rho_t+\diverg(\rho_t v_t)=0$ gives the weighted Poisson equation in~\eqref{eq:least-square-field-explicit}. The inverse notation is just a shorthand for solving this equation on zero-mean right-hand sides, modulo additive constants.
\end{proof}
\subsection{Benamou--Brenier dynamic formulation of OT}
\label{sec-benamou-brenier-dynamic}
\paragraph{Benamou--Brenier formulation.}
\begin{thm}[Benamou--Brenier]\label{thm-benamou-brenier}
For probability measures $\alpha_0,\alpha_1\in\Pp_2(\RR^d)$,
\begin{equation}\label{eq:benamou-brenier}
\Wass_2^2(\alpha_0,\alpha_1)
=
\inf_{(\alpha_t,v_t)}
\int_0^1\!\int_{\RR^d}\norm{v_t(x)}^2\,\d\alpha_t(x)\d t,
\end{equation}
where the infimum is over $(\alpha_t,v_t)$ solving $\partial_t\alpha_t+\nabla\!\cdot(\alpha_t v_t)=0$ with $\alpha_{t=0}=\alpha_0$ and $\alpha_{t=1}=\alpha_1$. If $\alpha_0$ has a density and $T$ is the optimal Monge map $T_\sharp\alpha_0=\alpha_1$, a minimizing curve is
\begin{equation}\label{eq:static-to-dynamic}
\alpha_t=((1-t)\Id+tT)_\sharp\alpha_0,
\qquad
v_t\big((1-t)x+tT(x)\big)=T(x)-x
\quad\text{for $\alpha_0$-a.e. $x$ and a.e. $t\in(0,1)$}.
\end{equation}
\end{thm}
\begin{proof}
For the inequality ``dynamic $\leq$ static'', assume first that a Monge map $T$ exists and define $(\alpha_t,v_t)$ by~\eqref{eq:static-to-dynamic}. Since the Lagrangian velocity $T(x)-x$ is independent of $t$,
\(\int_0^1\!\int\norm{v_t}^2\,\d\alpha_t\d t = \int\norm{T(x)-x}^2\,\d\alpha_0(x),\)
so the dynamic cost is no larger than the static Monge cost. Without a Monge map, the same construction is made with an optimal coupling $\pi$: sample $(X,Y)\sim\pi$ and move along the straight path $\gamma_{X,Y}(t)=(1-t)X+tY$. This path measure has action $\int\norm{x-y}^2\d\pi(x,y)$; projecting path velocities onto their conditional mean at time $t$ gives an admissible Eulerian velocity with no larger action, so the dynamic value is no larger than the Kantorovich value.

Conversely, for a smooth deterministic path, take the flow $T_t$ defined by $\dot T_t=v_t\circ T_t$ and $T_0=\Id$. Then $\alpha_t=(T_t)_\sharp\alpha_0$ and $(T_1)_\sharp\alpha_0=\alpha_1$. Jensen's inequality gives
\(\norm{T_1(x)-x}^2 \leq \int_0^1\norm{v_t(T_t(x))}^2\d t.\)
After integration with respect to $\alpha_0$, the Monge cost is bounded above by the dynamic action. For general finite-energy solutions of the continuity equation, the superposition principle lifts the curve to a probability measure on absolutely continuous paths; applying Jensen's inequality pathwise gives a coupling of the endpoints whose quadratic cost is no larger than the action. Thus the Kantorovich value is bounded above by the dynamic value.
\end{proof}
\paragraph{Convex moment-based reformulation.}
Although~\eqref{eq:benamou-brenier} is not jointly convex in $(\alpha_t,v_t)$, it becomes convex after replacing velocities by momenta. Given a velocity \(v\in L^2(\alpha;\RR^d)\), define the momentum

\(\omega\eqdef \alpha v, \qquad \omega(A)=\int_A v(x)\,\d\alpha(x),\)
\begin{equation}\label{eq-quadratic-perspective}
J(a,m)
\eqdef
\begin{cases}
\norm{m}^2/a, & a>0,\\
0, & a=0\ \text{and}\ m=0,\\
+\infty, & a=0\ \text{and}\ m\neq0,
\end{cases}
\qquad
(a,m)\in[0,+\infty)\times\RR^d.
\end{equation}
\begin{equation}\label{eq-measure-perspective-action}
\mathbb J(\alpha,\omega)
\eqdef
\int
J\left(
\frac{\d\alpha}{\d\lambda}(x),
\frac{\d \omega}{\d\lambda}(x)
\right)\d\lambda(x).
\end{equation}
\[
\mathbb J(\alpha,\omega)<+\infty
\quad\Longleftrightarrow\quad
\omega=v\alpha\ \text{with}\ v\in L^2(\alpha;\RR^d),
\qquad
\mathbb J(\alpha,\omega)=\int\norm{v}^2\,\d\alpha.
\]
\begin{equation}\label{eq:benamou-brenier-convex}
\Wass_2^2(\alpha_0,\alpha_1)
=
\inf_{\substack{\partial_t\alpha_t+\diverg \omega_t=0\\
\alpha_{t=0}=\alpha_0,\ \alpha_{t=1}=\alpha_1}}
\int_0^1 \mathbb J(\alpha_t,\omega_t)\,\d t.
\end{equation}
\[
\inf_{\substack{\partial_t\rho_t+\diverg m_t=0\\
\rho_{t=0}\d x=\alpha_0,\ \rho_{t=1}\d x=\alpha_1}}
\int_0^1\!\int_{\RR^d}
J(\rho_t(x),m_t(x))\,\d x\,\d t
=
\inf_{\substack{\partial_t\rho_t+\diverg m_t=0\\
\rho_{t=0}\d x=\alpha_0,\ \rho_{t=1}\d x=\alpha_1}}
\int_0^1\!\int_{\RR^d}
\frac{\norm{m_t(x)}^2}{\rho_t(x)}\,\d x\,\d t,
\]
\paragraph{Dual Hamilton--Jacobi formulation.}
The momentum formulation also has a useful dual. It turns the least-action problem into a Hamilton--Jacobi subsolution inequality for a scalar potential, with equality on the part of space-time actually visited by the optimal curve.

\begin{prop}[Dual Benamou--Brenier problem]\label{prop-benamou-brenier-dual}
Assume, for simplicity, that the densities are smooth, compactly supported, and that boundary terms vanish. Then the convex dynamic value~\eqref{eq:benamou-brenier-convex} has the dual formulation
\begin{equation}\label{eq:benamou-brenier-dual}
\Wass_2^2(\alpha_0,\alpha_1)
=
\sup_{\phi}
\enscond{
\int_{\RR^d}\phi_1\,\d\alpha_1
-
\int_{\RR^d}\phi_0\,\d\alpha_0
}{
\partial_t\phi_t+\frac14\norm{\nabla\phi_t}^2\leq 0
},
\end{equation}
where the supremum is over smooth test potentials $\phi:[0,1]\times\RR^d\to\RR$. In the general metric-measure statement, the same formula is understood after relaxation: $\phi$ is a Hamilton--Jacobi subsolution, for instance in the viscosity sense, and finite-dimensional discretizations follow directly from Fenchel--Rockafellar duality.
If $(\rho,m)$ and $\phi$ are smooth primal and dual optimizers, then
\begin{equation}\label{eq:bb-primal-dual-relation}
m_t=\frac{\rho_t}{2}\nabla\phi_t,
\qquad
\partial_t\phi_t+\frac14\norm{\nabla\phi_t}^2=0
\quad\text{on }\{\rho_t>0\}.
\end{equation}
Equivalently, the optimal Eulerian velocity is $v_t=m_t/\rho_t=\nabla\phi_t/2$.
\end{prop}
\begin{proof}
Let $(\rho,m)$ satisfy the continuity equation and let $\phi$ be smooth. Multiplying the constraint by $\phi$ and integrating by parts gives
\[
\int_{\RR^d}\phi_1\,\d\alpha_1
-
\int_{\RR^d}\phi_0\,\d\alpha_0
=
\int_0^1\!\int_{\RR^d}
\left(\rho_t\,\partial_t\phi_t+\dotp{m_t}{\nabla\phi_t}\right)\d x\d t.
\]
If $\partial_t\phi_t+\norm{\nabla\phi_t}^2/4\leq0$, Young's inequality yields, pointwise, \(\rho\,\partial_t\phi+\dotp{m}{\nabla\phi} \leq -\frac{\rho}{4}\norm{\nabla\phi}^2+\dotp{m}{\nabla\phi} \leq J(\rho,m),\)
with the perspective convention~\eqref{eq-quadratic-perspective}. Hence the dual objective of every feasible $\phi$ is no larger than the action of every feasible primal pair.

Conversely, introduce $\phi$ as a Lagrange multiplier for $\partial_t\rho+\diverg m=0$. After integration by parts, and up to the fixed endpoint contribution in~\eqref{eq:benamou-brenier-dual}, the pointwise minimization over $(\rho,m)$ contains
\(J(\rho,m)-\rho\,\partial_t\phi-\dotp{m}{\nabla\phi}.\)
For fixed $\rho>0$, its minimum over $m$ is attained at $m=\rho\nabla\phi/2$ and equals $-\rho(\partial_t\phi+\norm{\nabla\phi}^2/4)$. The recession cases in~\eqref{eq-quadratic-perspective} give the same conclusion when \(\rho=0\). Thus the infimum over $\rho\geq0$ is finite exactly when the Hamilton--Jacobi inequality holds, in which case the pointwise infimum is $0$. Fenchel--Rockafellar duality then gives no duality gap. Equality in the two pointwise inequalities gives~\eqref{eq:bb-primal-dual-relation}.
\end{proof}
\(\frac{\d}{\d t}\phi_t(\gamma(t)) = \partial_t\phi_t(\gamma(t))+\dotp{\nabla\phi_t(\gamma(t))}{\dot\gamma(t)} \leq \norm{\dot\gamma(t)}^2,\)
where the last inequality uses $\partial_t\phi+\norm{\nabla\phi}^2/4\leq0$ and Young's inequality. Integrating and minimizing over curves gives

\(\phi_1(y)-\phi_0(x)\leq\norm{x-y}^2.\)
Thus $(-\phi_0,\phi_1)$ is a feasible static Kantorovich dual pair for the quadratic cost. At optimality the inequality is saturated on the endpoint pairs connected by the primal characteristics, which is the Eulerian version of complementary slackness.

\paragraph{Proximal splitting.}
The convex momentum formulation also explains the original Benamou--Brenier solver. Suppressing discretization indices, write \(U=(\rho,m)\), where \(m\) is the density of the momentum measure.

\[
\mathcal F(U)=\int_0^1\!\int_{\RR^d}J(\rho_t(x),m_t(x))\,\d x\,\d t,
\qquad
\mathcal C=\enscond{(\rho,m)}{\partial_t\rho+\diverg m=0,\ \rho_0\d x=\alpha_0,\ \rho_1\d x=\alpha_1},
\]
\(\min_U\; \mathcal F(U)+\mathcal G(U).\)
\[
\prox_{\tau\mathcal F}(\bar U)
=
\argmin_U
\frac12\norm{U-\bar U}_{\mathsf H}^2+\tau\mathcal F(U),
\qquad
\prox_{\tau\mathcal G}(\bar U)
=
\argmin_{U\in\mathcal C}
\frac12\norm{U-\bar U}_{\mathsf H}^2.
\]
For the standard product metric, the first map is local in $(t,x)$: it is the proximal operator of the convex perspective $J$. The second map is the orthogonal projection onto the affine set defined by the divergence equation and endpoint constraints. Douglas--Rachford then alternates these two simple operations:

\[
\begin{aligned}
U^{k+1} &= \prox_{\tau\mathcal F}(Z^k),\\
\widetilde U^{k+1} &= \prox_{\tau\mathcal G}(2U^{k+1}-Z^k),\\
Z^{k+1} &= Z^k+\widetilde U^{k+1}-U^{k+1}.
\end{aligned}
\]
\paragraph{Path-space formulation.}\label{rem-bb-path-space}
Let $\Ss=C([0,1];\RR^d)$ be the space of continuous paths endowed with the uniform topology. For $t\in[0,1]$ define the evaluation map

\(e_t:\Ss\to\RR^d, \qquad e_t(\gamma)=\gamma(t).\)
\[
\Wass_2^2(\alpha_0,\alpha_1)
=
\inf_{M\in\Pp(\Ss)}
\enscond{
\int_{\Ss}\!\int_0^1\norm{\dot\gamma(t)}^2\d t\,\d M(\gamma)
}{
(e_0)_\sharp M=\alpha_0,\ (e_1)_\sharp M=\alpha_1
}.
\]
\((e_t,\dot e_t)_\sharp M^*(\d x,\d q) = \alpha_t(\d x)\delta_{v_t^*(x)}(\d q),\)
where $v_t^*$ is the optimal velocity field in the Benamou--Brenier formulation. Hence $M^*$ concentrates on straight-line geodesics and, for a.e. $t$, assigns exactly one direction at $\alpha_t$-a.e. spatial point.

\subsection{Generalized Dynamic Wasserstein Distances}
\label{sec-bb-extensions}
\label{sec-generalized-dynamic-wasserstein-distances}
The quadratic Benamou--Brenier formula is only one instance of a broader fixed-mass dynamic language. The objects introduced here are metric: they specify admissible curves, tangent variables and path energies.

\paragraph{Path actions.}
In the mass-preserving Euclidean setting, the basic input is an instantaneous action $\mathbb A(\alpha,w)$, where \(\alpha\) is the current measure and \(w\) is an admissible velocity representative. When this action is normalized as a squared infinitesimal speed, it generates the length-space value

\begin{equation}\label{eq-generalized-action-length-distance}
\mathsf D_{\mathbb A}^2(\alpha_0,\alpha_1)
=
\inf_{\alpha_t,v_t}
\left\{
\int_0^1 \mathbb A(\alpha_t,v_t)\,\d t
:
\partial_t\alpha_t+
\diverg(\alpha_t v_t)=0,\
\alpha_{t=0}=\alpha_0,\
\alpha_{t=1}=\alpha_1
\right\}.
\end{equation}
\paragraph{Quadratic, or Riemannian, tangent actions.}
\[
\mathbb A(\alpha;w,z)
=
\left\langle Q_\alpha w,z\right\rangle_{L^2(\alpha)},
\qquad
\mathbb A(\alpha,w)=\left\langle Q_\alpha w,w\right\rangle_{L^2(\alpha)},
\]
\begin{equation}\label{eq-general-quadratic-tangent-action}
\mathsf D_Q^2(\alpha_0,\alpha_1)
=
\mathsf D_{\mathbb A}^2(\alpha_0,\alpha_1)
=
\inf_{\substack{\partial_t\alpha_t+\diverg(\alpha_t v_t)=0\\
\alpha_{t=0}=\alpha_0,\ \alpha_{t=1}=\alpha_1}}
\int_0^1
\left\langle Q_{\alpha_t}v_t,v_t\right\rangle_{L^2(\alpha_t)}
\d t.
\end{equation}
The usual $\Wass_2$ geometry corresponds to $Q_\alpha=\Id$ in this simplified notation. Thus \(Q_\alpha\) records how the chosen geometry deforms the Euclidean \(L^2(\alpha)\) tangent norm: no deformation for \(\Wass_2\), and a nontrivial tensor for generalized Riemannian geometries.

\paragraph{Local velocity actions.}
\(A:[0,+\infty)\times\RR^d\to[0,+\infty], \qquad (a,w)\mapsto A(a,w),\)
where $a\in\RR_+$ is a density value and $w\in\RR^d$ is a velocity value, and defines

\begin{equation}\label{eq-local-velocity-action}
\mathbb A(\alpha,w)
=
\int A\left(\frac{\d\alpha}{\d\lambda}(x),w(x)\right)\d\lambda(x).
\end{equation}
For a fixed reference $\lambda$, this covers density-dependent mobilities and congestion constraints. If, in addition, $A$ is positively $1$-homogeneous in its first variable, $A(\eta a,w)=\eta A(a,w)$ for $\eta\geq0$, then the same formula is intrinsic: replacing $\lambda$ by another dominating measure gives the same value. The usual Benamou--Brenier action is the model case

\begin{equation}\label{eq-bb-velocity-action}
A_2(a,w)=a\norm{w}^2,
\qquad
\mathbb A(\alpha,w)=\int\norm{w}^2\d\alpha.
\end{equation}
\paragraph{Homogeneous momentum actions.}\label{rem-generalized-bb}
\begin{equation}\label{eq-general-momentum-perspective}
J_A(a,m)
\eqdef
\begin{cases}
A(a,m/a), & a>0,\\
0, & a=0\ \text{and}\ m=0,\\
+\infty, & a=0\ \text{and}\ m\neq0,
\end{cases}
\end{equation}
\begin{equation}\label{eq-general-measure-momentum-action}
\mathbb J_{A,\lambda}(\alpha,\omega)
\eqdef
\int
J_A\!\left(
\frac{\d\alpha}{\d\lambda},
\frac{\d\omega}{\d\lambda}
\right)\d\lambda,
\end{equation}
\(J_2(a,m)=\frac{\norm m^2}{a},\)
\begin{prop}[Concave mobilities give convex momentum actions]\label{prop-momentum-perspective-convexity}
Let \(I\subset[0,+\infty)\) be a convex interval, let \(\theta:I\to[0,+\infty)\) be concave, and let \(L:\RR^d\to[0,+\infty]\) be convex with \(L(0)=0\). Define, on the set where \(\theta(a)>0\),
\begin{equation}\label{eq-concave-mobility-perspective}
J_{\theta,L}(a,m)
\eqdef
\theta(a)L\!\left(\frac{m}{\theta(a)}\right),
\qquad
A_{\theta,L}(a,w)
\eqdef
J_{\theta,L}(a,aw)
=
\theta(a)L\!\left(\frac{aw}{\theta(a)}\right),
\end{equation}
and extend \(J_{\theta,L}\) to the boundary by lower semicontinuity. Then \(J_{\theta,L}\) is convex in \((a,m)\). In particular, this single construction contains
\(A(a,w)=aL(w) \quad\text{by taking }\theta(a)=a,\)
and the concave-mobility quadratic action
\(A(a,w)=\frac{a^2\norm{w}^2}{\theta(a)} \quad\text{by taking }L(u)=\norm u^2.\)
\end{prop}
\begin{proof}
The perspective \(P(s,m)=sL(m/s)\) of a convex function is convex on \(s>0\). Since \(L(0)=0\), it is also nonincreasing in the scalar \(s\) for fixed \(m\): if \(s_2\geq s_1>0\), then
\(L(m/s_2) = L\big((s_1/s_2)(m/s_1)\big) \leq (s_1/s_2)L(m/s_1),\)
hence \(P(s_2,m)\leq P(s_1,m)\). Let \(a_\zeta=(1-\zeta)a_0+\zeta a_1\) and \(m_\zeta=(1-\zeta)m_0+\zeta m_1\). By concavity of \(\theta\), \(\theta(a_\zeta) \geq (1-\zeta)\theta(a_0)+\zeta\theta(a_1).\)
Using the monotonicity in \(s\), and then the convexity of \(P\), gives
\[
J_{\theta,L}(a_\zeta,m_\zeta)
=
P(\theta(a_\zeta),m_\zeta)
\leq
P((1-\zeta)\theta(a_0)+\zeta\theta(a_1),m_\zeta)
\leq
(1-\zeta)J_{\theta,L}(a_0,m_0)+\zeta J_{\theta,L}(a_1,m_1).
\]
The boundary cases follow by the chosen lower-semicontinuous extension.
\end{proof}
\begin{prop}[Homogeneous dynamic actions define distances]\label{prop-homogeneous-dynamic-action-distance}
Fix a reference measure \(\lambda\), omitted from the notation only in the intrinsic case where \(\mathbb J_{A,\lambda}\) does not depend on \(\lambda\). Assume that the momentum perspective \(J_A\) defined in~\eqref{eq-general-momentum-perspective} is lower semicontinuous, convex in \((a,m)\), even in \(m\), and satisfies \(J_A(a,0)=0\). Assume moreover that, for some \(r>1\),
\(J_A(a,\xi m)=|\xi|^rJ_A(a,m) \qquad \text{for all admissible }a,\ \text{all }m,\ \text{and all }\xi\in\RR,\)
and that \(J_A(a,m)=0\) if and only if \(m=0\). Equivalently, the evenness, homogeneity and nondegeneracy assumptions translate, away from zero density, into
\(A(a,-w)=A(a,w), \qquad A(a,\xi w)=|\xi|^rA(a,w), \qquad A(a,w)=0 \Longleftrightarrow w=0 \qquad (a>0).\)
Define, on every fixed-mass class,
\[
\mathsf D_{A,\lambda}(\alpha_0,\alpha_1)
\eqdef
\inf_{\substack{\partial_t\alpha_t+\diverg\omega_t=0\\
\alpha_{t=0}=\alpha_0,\ \alpha_{t=1}=\alpha_1}}
\left(
\int_0^1\mathbb J_{A,\lambda}(\alpha_t,\omega_t)\,\d t
\right)^{1/r}.
\]
Assume finally that the relaxed dynamic problem is sequentially closed and attains its infimum whenever the value is finite. Then \(\mathsf D_{A,\lambda}\) is an extended distance on each finite-action component: it is symmetric, satisfies the triangle inequality, and \(\mathsf D_{A,\lambda}(\alpha_0,\alpha_1)=0\) only when \(\alpha_0=\alpha_1\). In the intrinsic case, we write \(\mathsf D_A\).
\end{prop}
\begin{proof}
The constant curve \(\alpha_t=\alpha_0\), \(\omega_t=0\), gives \(\mathsf D_{A,\lambda}(\alpha_0,\alpha_0)=0\). Conversely, if \(\mathsf D_{A,\lambda}(\alpha_0,\alpha_1)=0\), attainment provides a relaxed zero-action minimizer, which satisfies
\(\int_0^1 \mathbb J_{A,\lambda}(\alpha_t,\omega_t)\,\d t=0.\)
Since \(\mathbb J_{A,\lambda}\geq0\), this implies \(\omega_t=0\) for a.e. \(t\). The continuity equation then gives \(\partial_t\alpha_t=0\) in the sense of distributions, so \(\alpha_0=\alpha_1\).

Symmetry follows by time reversal. If \((\alpha_t,\omega_t)_{t\in[0,1]}\) connects \(\alpha_0\) to \(\alpha_1\), then
\(\widetilde\alpha_t=\alpha_{1-t}, \qquad \widetilde\omega_t=-\omega_{1-t}\)
connects \(\alpha_1\) to \(\alpha_0\), satisfies the same continuity equation, and has the same action because \(J_A\) is even in \(m\).

Let \((\alpha^1,\omega^1)\) connect \(\alpha_0\) to \(\alpha_1\) with action \(E_1\), and let \((\alpha^2,\omega^2)\) connect \(\alpha_1\) to \(\alpha_2\) with action \(E_2\). For \(\tau\in(0,1)\), concatenate the two curves after the time changes \(s=t/\tau\) and \(s=(t-\tau)/(1-\tau)\):
\[
(\alpha_t,\omega_t)=
\begin{cases}
\bigl(\alpha^1_{t/\tau},\,\tau^{-1}\omega^1_{t/\tau}\bigr), & 0\leq t\leq \tau,\\[.2em]
\bigl(\alpha^2_{(t-\tau)/(1-\tau)},\,(1-\tau)^{-1}\omega^2_{(t-\tau)/(1-\tau)}\bigr), & \tau\leq t\leq1.
\end{cases}
\]
The rescaled curve is admissible from \(\alpha_0\) to \(\alpha_2\). By \(r\)-homogeneity in the momentum variable, its action is \(\tau^{1-r}E_1+(1-\tau)^{1-r}E_2.\)
Writing \(X=E_1^{1/r}\) and \(Y=E_2^{1/r}\), the optimal allocation is \(\tau=X/(X+Y)\), with the evident limiting convention if \(X+Y=0\), and gives
\[
\inf_{\tau\in(0,1)}
\left\{\tau^{1-r}E_1+(1-\tau)^{1-r}E_2\right\}
=
(X+Y)^r.
\]
Hence \(\mathsf D_{A,\lambda}(\alpha_0,\alpha_2)\leq E_1^{1/r}+E_2^{1/r}\). Taking \(E_1\) and \(E_2\) arbitrarily close to the two optimal action values yields
\(\mathsf D_{A,\lambda}(\alpha_0,\alpha_2) \leq \mathsf D_{A,\lambda}(\alpha_0,\alpha_1)+\mathsf D_{A,\lambda}(\alpha_1,\alpha_2),\)
with the usual extended-real convention when one of the distances is infinite.
\end{proof}
\paragraph{Concave-mobility actions.}
\[
J_\theta(a,m)
\eqdef
\begin{cases}
\norm{m}^2/\theta(a), & \theta(a)>0,\\
0, & \theta(a)=0 \text{ and } m=0,\\
+\infty, & \theta(a)=0 \text{ and } m\neq0.
\end{cases}
\]
\(A_\theta(a,w)=J_\theta(a,aw)=\frac{a^2\norm{w}^2}{\theta(a)}.\)
\(\mathbb A_{\theta,\lambda}(\alpha,w) = \int A_\theta(a(x),w(x))\,\d\lambda(x) = \int \frac{a(x)^2}{\theta(a(x))}\norm{w(x)}^2\,\d\lambda(x),\)
\(\mathbb A_{\theta,\lambda}(\alpha,w) = \int \frac{a(x)}{\theta(a(x))}\norm{w(x)}^2\,\d\alpha(x).\)
Hence this is a local Riemannian case, in the sense of~\eqref{eq-general-quadratic-tangent-action}, whenever the multiplier \(a(x)/\theta(a(x))\) is finite and positive. The associated tensor is the multiplication operator

\begin{equation}\label{eq-concave-mobility-q-alpha}
\big(Q_{\theta,\lambda,\alpha}w\big)(x)
=
\frac{a(x)}{\theta(a(x))}w(x),
\qquad
\alpha=a\lambda,
\end{equation}
defined \(\alpha\)-a.e. on the set where \(\theta(a)>0\). Except for linear mobilities \(\theta(a)=ca\), and in particular the normalized case \(\theta(a)=a\) which recovers \(\Wass_2\), the pointwise velocity action \(A_\theta(a,w)\) is not positively \(1\)-homogeneous in the density variable \(a\).

\(\mathsf W_{\theta,\lambda}(\alpha_0,\alpha_1) \eqdef \mathsf D_{A_\theta,\lambda}(\alpha_0,\alpha_1),\)
where the subscript \(\lambda\) recalls that the action is measured through the density \(a=\d\alpha/\d\lambda\). Equivalently, because this action is quadratic in the momentum, \(\mathsf W_{\theta,\lambda}^2\) is the path value~\eqref{eq-generalized-action-length-distance} with \(\mathbb A=\mathbb A_{\theta,\lambda}\).

\begin{prop}[Concave-mobility dynamic distances]\label{prop-concave-mobility-distance}
For a fixed reference measure \(\lambda\), and under the standard compactness and lower-semicontinuity hypotheses ensuring existence of relaxed minimizers, \(\mathsf W_{\theta,\lambda}=\mathsf D_{A_\theta,\lambda}\) is an extended distance on each finite-action component contained in \(\{\alpha:\alpha\ll\lambda\}\). In particular, whenever the relevant component has finite pairwise distances, \(\mathsf W_{\theta,\lambda}\) is a genuine metric there.
\end{prop}
\begin{proof}
Proposition~\ref{prop-momentum-perspective-convexity}, applied with \(L(u)=\norm u^2\), gives the lower-semicontinuous convex momentum density \(J_\theta\). It is even in \(m\), satisfies \(J_\theta(a,0)=0\), and obeys
\(J_\theta(a,\xi m)=|\xi|^2J_\theta(a,m).\)
Moreover \(J_\theta(a,m)=0\) if and only if \(m=0\), with the boundary convention used in its definition. For the fixed reference \(\lambda\), the hypotheses of Proposition~\ref{prop-homogeneous-dynamic-action-distance} therefore hold with \(r=2\); the existence assumption is supplied by the compactness hypothesis in the statement. That proposition gives symmetry, separation and the triangle inequality for \(\mathsf D_{A_\theta,\lambda}\), hence for \(\mathsf W_{\theta,\lambda}\).
\end{proof}
\paragraph{Dynamic spectral Wasserstein distances.}
Let $\gamma$ be a monotone spectral gauge on $\mathbb S_+^d$. For a probability measure $\alpha$ and a velocity field $v\in L^2(\alpha;\RR^d)$, define the spectral tangent action

\begin{equation}\label{eq-spectral-tangent-action}
\mathbb A_\gamma(\alpha,v)
\eqdef
\gamma\!\left(\int v(x)v(x)^\top\d\alpha(x)\right).
\end{equation}
\begin{equation}\label{eq-dynamic-spectral-wasserstein}
\mathsf W_{\gamma,\mathrm{dyn}}^2
\eqdef
\mathsf D_{\mathbb A_\gamma}^2.
\end{equation}
\[
\mathbb J_\gamma(\alpha,\omega)
=
\gamma\!\left(\int
\left(\frac{\d\omega}{\d\alpha}\right)
\left(\frac{\d\omega}{\d\alpha}\right)^\top
\d\alpha\right),
\]
\[
\mathbb J_\gamma(\rho,m)
=
\gamma\!\left(\int \frac{m(x)m(x)^\top}{\rho(x)}\,\d x\right).
\]
\(A_{\mathrm{lin}}(a,w)=a\,w^\top G w, \qquad J_{\mathrm{lin}}(a,m)=\frac{m^\top G m}{a},\)
\begin{prop}[Static/dynamic equality for spectral Wasserstein]\label{prop-static-dynamic-spectral-wasserstein}
Let $\gamma$ be a monotone spectral gauge and let $\alpha_0,\alpha_1\in\Pp_2(\RR^d)$. Then
\(\mathsf W_{\gamma,\mathrm{dyn}}^2(\alpha_0,\alpha_1) = \Wass_\gamma^2(\alpha_0,\alpha_1).\)
The common value defines a distance on $\Pp_2(\RR^d)$; the triangle inequality is established by the robust representation in Proposition~\ref{prop-spectral-wasserstein-robust}.
\end{prop}
\begin{proof}
We first prove $\mathsf W_{\gamma,\mathrm{dyn}}^2\leq\Wass_\gamma^2$. Let $\pi\in\Couplings(\alpha_0,\alpha_1)$ and let $(X,Y)\sim\pi$. Set $Z_t=(1-t)X+tY$ and $\alpha_t=(Z_t)_\sharp\pi$. Define the Eulerian velocity as the conditional mean
\(v_t(z)\eqdef \EE[Y-X\mid Z_t=z].\)
Then $(\alpha_t,v_t)$ solves the continuity equation. If
\(M_\pi\eqdef\int(x-y)(x-y)^\top\d\pi(x,y), \qquad C_t\eqdef\int v_t(z)v_t(z)^\top\d\alpha_t(z),\)
conditional Jensen gives \(C_t\preceq M_\pi\). Indeed, for every direction \(u\in\RR^d\),
\[
u^\top C_tu
=
\EE\!\left[\EE[\dotp{u}{Y-X}\mid Z_t]^2\right]
\leq
\EE[(\dotp{u}{Y-X})^2]
=
u^\top M_\pi u.
\]
By the Loewner monotonicity of $\gamma$, \(\int_0^1\gamma(C_t)\d t \leq \gamma(M_\pi).\)
Taking the infimum over $\pi$ gives the first inequality.

Conversely, let $(\alpha_t,v_t)$ be an admissible finite-action competitor. Since $\gamma$ is a finite positive gauge on the finite-dimensional cone $\mathbb S_+^d$, it is equivalent to the trace on this cone; hence finite spectral action implies the usual finite kinetic energy needed for the superposition principle. Thus there exists a probability law $\eta$ on absolutely continuous paths $\omega_\cdot$ such that $(e_t)_\sharp\eta=\alpha_t$ and $\dot\omega_t=v_t(\omega_t)$ for a.e. \(t\). Let $\pi=(e_0,e_1)_\sharp\eta$ be the induced endpoint coupling. With \(C_t=\int v_tv_t^\top\d\alpha_t\), Jensen's inequality along each path gives
\[
(\omega_1-\omega_0)(\omega_1-\omega_0)^\top
=
\left(\int_0^1\dot\omega_t\d t\right)
\left(\int_0^1\dot\omega_t\d t\right)^\top
\preceq
\int_0^1\dot\omega_t\dot\omega_t^\top\d t.
\]
After integration over $\eta$, \(M_\pi\preceq \int_0^1 C_t\d t.\)
Monotonicity of $\gamma$, followed by convexity, yields
\[
\gamma(M_\pi)
\leq
\gamma\!\left(\int_0^1 C_t\d t\right)
\leq
\int_0^1\gamma(C_t)\d t.
\]
Since $\pi$ is admissible for the static problem, $\Wass_\gamma^2(\alpha_0,\alpha_1)$ is bounded by the action of the dynamic competitor. Taking the infimum over all competitors gives the reverse inequality.

The crucial assumption is the monotonicity of $\gamma$: both parts of the proof produce Loewner-order comparisons between covariance matrices, and these comparisons imply inequalities between actions only for monotone gauges.
\end{proof}
\paragraph{Kernelized Benamou--Brenier distances.}\label{sec-kernelized-bb-distance}
Let $k$ be a positive definite kernel on $\RR^d$ with scalar RKHS $\RKHS_k$. The vector-valued RKHS is

\(\RKHS_k^d\eqdef \RKHS_k\times\cdots\times\RKHS_k, \qquad \norm{v}_{\RKHS_k^d}^2 \eqdef \sum_{\ell=1}^d\norm{v_\ell}_{\RKHS_k}^2.\)
\begin{equation}\label{eq-kernelized-bb-distance}
\mathbb A_k(\alpha,v)
\eqdef
\norm{v}_{\RKHS_k^d}^2,
\qquad
\mathcal W_k^2
\eqdef
\mathsf D_{\mathbb A_k}^2,
\end{equation}
where the general distance formula~\eqref{eq-generalized-action-length-distance} is understood with the restricted admissible tangent class \(v_t\in\RKHS_k^d\). The action itself is independent of \(\alpha\); the measure only enters through the continuity equation \(\partial_t\alpha_t+\diverg(\alpha_t v_t)=0\), which says how the common smooth velocity field moves all particles.

\begin{prop}[Kernelized dynamic distance]\label{prop-kernelized-bb-distance}
Assume that the kernel has uniformly bounded diagonal, \(\kappa_k^2\eqdef\sup_{x\in\RR^d}k(x,x)<+\infty.\)
Then the quantity \(\mathcal W_k\) defined by~\eqref{eq-kernelized-bb-distance} is an extended distance on probability measures. More precisely, it is symmetric, satisfies the triangle inequality, and \(\mathcal W_k(\alpha_0,\alpha_1)=0\) implies \(\alpha_0=\alpha_1\). It can take the value \(+\infty\) between different finite-action components.
\end{prop}
\begin{proof}
The constant curve gives zero self-distance. To prove separation without assuming existence of a minimizing curve, use the RKHS evaluation bound
\(\norm{v(x)}\leq \kappa_k\norm{v}_{\RKHS_k^d}.\)
For every $\varphi\in C_c^1(\RR^d)$ and every admissible curve of action $E$, the weak continuity equation and Cauchy--Schwarz give
\[
\begin{aligned}
\left|\int\varphi\,\d(\alpha_1-\alpha_0)\right|
&\leq
\int_0^1\!\int\norm{\nabla\varphi(x)}\,\norm{v_t(x)}\,\d\alpha_t(x)\d t\\
&\leq
\kappa_k\norm{\nabla\varphi}_\infty
\left(\int_0^1\norm{v_t}_{\RKHS_k^d}^2\d t\right)^{1/2}
=
\kappa_k\norm{\nabla\varphi}_\infty\sqrt E.
\end{aligned}
\]
Taking the infimum over admissible curves shows that $\mathcal W_k(\alpha_0,\alpha_1)=0$ forces equality of the integrals of every compactly supported smooth test function, hence $\alpha_0=\alpha_1$.

Symmetry follows by time reversal and by replacing \(v_t\) with \(-v_{1-t}\). For the triangle inequality, concatenate an almost optimal curve from \(\alpha_0\) to \(\alpha_1\) of action \(E_1\) with one from \(\alpha_1\) to \(\alpha_2\) of action \(E_2\). Compressing the first curve into a time interval of length \(\tau\) multiplies its action by \(1/\tau\), and compressing the second into length \(1-\tau\) multiplies its action by \(1/(1-\tau)\). Optimizing over \(\tau\in(0,1)\) gives \((\sqrt{E_1}+\sqrt{E_2})^2\), and taking infima yields the triangle inequality.
\end{proof}
One should read \(\mathcal W_k\) as an extended distance on finite-action components, not as a replacement for $\Wass_2$ on all of $\Pp_2(\RR^d)$.

\subsection{Nonlocal Wasserstein Distances}\label{sec-nonlocal-wasserstein-distances}
\begin{equation}\label{eq-logarithmic-mean}
\theta(a,b)\eqdef
\begin{cases}
\displaystyle\frac{a-b}{\log a-\log b}, & a\neq b,\\[.4em]
a, & a=b,
\end{cases}
\end{equation}
\paragraph{Continuum jump kernels.}
\begin{equation}\label{eq-nonlocal-jump-measure}
\int_{\mathcal X\times\mathcal X}\Phi(x,y)\,\mathsf J(\d x,\d y)
\eqdef
\int_{\mathcal X}\left(\int_{\mathcal X}\Phi(x,y)\,K(x,\d y)\right)\mathfrak m(\d x).
\end{equation}
\(\bar\nabla \varphi(x,y)\eqdef \varphi(y)-\varphi(x)\)
for the nonlocal gradient, and use the logarithmic mean \(\theta\) defined in~\eqref{eq-logarithmic-mean}. A curve \(\alpha_t=\rho_t\mathfrak m\) driven by an antisymmetric velocity \(v_t(x,y)=-v_t(y,x)\) satisfies the nonlocal continuity equation if, for all test functions \(\varphi\),

\begin{equation}\label{eq-nonlocal-continuity-weak}
\frac{\d}{\d t}\int \varphi\,\d\alpha_t
=
\frac12
\iint
\bar\nabla\varphi(x,y)\,
v_t(x,y)\,
\theta(\rho_t(x),\rho_t(y))\,
\mathsf J(\d x,\d y).
\end{equation}
\(\mathbb A_K(\alpha,v) \eqdef \frac12 \iint |v(x,y)|^2 \theta(\rho(x),\rho(y))\, \mathsf J(\d x,\d y),\)
\begin{equation}\label{eq-nonlocal-wasserstein-distance}
\mathcal W_K^2(\alpha_0,\alpha_1)
\eqdef
\inf_{\rho_t,v_t}
\int_0^1
\mathbb A_K(\alpha_t,v_t)\,\d t,
\end{equation}
where the infimum is over all curves solving \eqref{eq-nonlocal-continuity-weak} with endpoints \(\alpha_0,\alpha_1\).

\begin{prop}[Metric properties of nonlocal transport]\label{prop-nonlocal-distance-properties}
Under standard regularity and irreducibility assumptions on the jump kernel, \(\mathcal W_K\) is an extended distance on the set of probability measures absolutely continuous with respect to \(\mathfrak m\), and finite-distance pairs are connected by constant-speed geodesics.
\end{prop}
\begin{proof}
We use the analytic compactness and lower-semicontinuity theorem for the logarithmic-mean action. In particular, action-bounded sequences of admissible curves are compact for the narrow topology, the weak nonlocal continuity equation is closed under this convergence, and the action is lower semicontinuous. These facts are the nonlocal analogues of the compactness theorem used in the Benamou--Brenier formulation.

Nonnegativity is immediate from the definition of \(\mathbb A_K(\alpha,v)\). If \(\alpha_0=\alpha_1\), the constant curve \(\rho_t=\rho_0\), \(v_t=0\), is admissible and has zero action, hence \(\mathcal W_K(\alpha_0,\alpha_0)=0\).

Symmetry follows by time reversal. If \((\rho_t,v_t)\) transports \(\alpha_0\) to \(\alpha_1\), set
\(\tilde\rho_t=\rho_{1-t}, \qquad \tilde v_t=-v_{1-t}.\)
For every test function \(\varphi\), differentiating \(\int\varphi\tilde\rho_t\,\d\mathfrak m\) and using \eqref{eq-nonlocal-continuity-weak} shows that \((\tilde\rho_t,\tilde v_t)\) solves the same continuity equation from \(\alpha_1\) to \(\alpha_0\). Since the action is quadratic in \(v\), \(\mathbb A_K(\tilde\alpha_t,\tilde v_t)=\mathbb A_K(\alpha_{1-t},v_{1-t})\), so the action is unchanged.

For the triangle inequality, let \((\rho^0_t,v^0_t)\) connect \(\alpha_0\) to \(\alpha_1\) with action \(E_0\), and let \((\rho^1_t,v^1_t)\) connect \(\alpha_1\) to \(\alpha_2\) with action \(E_1\). For \(0<\zeta<1\), concatenate the two curves after the rescaling
\[
(\rho_t,v_t)=
\begin{cases}
\bigl(\rho^0_{t/\zeta},\,\zeta^{-1}v^0_{t/\zeta}\bigr),
&0\leq t\leq\zeta,\\[.35em]
\bigl(\rho^1_{(t-\zeta)/(1-\zeta)},\,(1-\zeta)^{-1}v^1_{(t-\zeta)/(1-\zeta)}\bigr),
&\zeta<t\leq1.
\end{cases}
\]
The factors \(\zeta^{-1}\) and \((1-\zeta)^{-1}\) are exactly those needed for the weak continuity equation after changing time. Since \(v\mapsto\mathbb A_K(\alpha,v)\) is quadratic,
\(\int_0^1\mathbb A_K(\alpha_t,v_t)\,\d t = \frac{E_0}{\zeta}+\frac{E_1}{1-\zeta}.\)
Optimizing in \(\zeta\), or taking \(\zeta=\sqrt{E_0}/(\sqrt{E_0}+\sqrt{E_1})\) when both actions are positive, gives a concatenated action \((\sqrt{E_0}+\sqrt{E_1})^2\). Taking infima over the two curves yields
\(\mathcal W_K(\alpha_0,\alpha_2) \leq \mathcal W_K(\alpha_0,\alpha_1) + \mathcal W_K(\alpha_1,\alpha_2).\)

If \(\mathcal W_K(\alpha_0,\alpha_1)=0\), choose admissible curves with actions tending to zero. By the compactness and lower-semicontinuity theorem quoted above, a subsequence converges to an admissible curve \((\rho_t,v_t)\) of zero action. Hence \(v_t=0\) for \(\theta(\rho_t(x),\rho_t(y))\mathsf J(\d x,\d y)\d t\)-a.e. \((t,x,y)\). Substituting in the weak continuity equation gives
\(\frac{\d}{\d t}\int\varphi\,\d\alpha_t=0\)
for every admissible test function \(\varphi\). The irreducibility/separation assumption in ensures that these test functions determine the measure, so \(\alpha_t\) is constant and \(\alpha_0=\alpha_1\).

Finally, if \(\mathcal W_K(\alpha_0,\alpha_1)<+\infty\), the same direct-method compactness applied to a minimizing sequence gives a minimizer of \eqref{eq-nonlocal-wasserstein-distance}. Reparametrizing this minimizing curve by metric arclength gives a constant-speed curve. Equivalently, for every \(0\leq s<t\leq1\), the restriction of the minimizer to \([s,t]\), rescaled to unit time, is admissible between \(\alpha_s\) and \(\alpha_t\); the triangle inequality and equality of the total minimal action force
\(\mathcal W_K(\alpha_s,\alpha_t)=(t-s)\mathcal W_K(\alpha_0,\alpha_1)\)
after this constant-speed parametrization. Thus finite-distance pairs are joined by constant-speed geodesics. The consequences for entropy dynamics and fractional PDE examples are developed in Section~\ref{sec-nonlocal-wasserstein-pdes}.
\end{proof}
For a fixed jump kernel, this geometry is genuinely nonlocal and does not coincide with ordinary $\Wass_2$. The local metric is nevertheless recovered in a small-jump limit. More explicitly, on $\mathcal X=\RR^d$, let $\eta(z)=\bar\eta(\norm{z})$ be a nonnegative radial profile with

\(M_2(\eta)\eqdef\int_{\RR^d}\norm{z}^2\eta(z)\,\d z\in(0,+\infty),\)
\begin{equation}\label{eq-small-jump-kernel-scaling}
K_\varepsilon(x,\d y)
\eqdef
\eta_\varepsilon(y-x)\,\d y,
\qquad
\eta_\varepsilon(z)
\eqdef
\varepsilon^{-d}\eta(z/\varepsilon).
\end{equation}
\begin{equation}\label{eq-small-jump-kernel-moments}
\int\norm{y-x}^2K_\varepsilon(x,\d y)
=
\varepsilon^2M_2(\eta),
\qquad
\int (y-x)(y-x)^\top K_\varepsilon(x,\d y)
=
\varepsilon^2\frac{M_2(\eta)}{d}\Id.
\end{equation}
\(\widehat K_\varepsilon \eqdef \frac{2d}{\varepsilon^2M_2(\eta)}K_\varepsilon, \qquad \int (y-x)(y-x)^\top\widehat K_\varepsilon(x,\d y)=2\Id.\)
\[
\mathcal W_{\widehat K_\varepsilon}
=
\varepsilon\sqrt{\frac{M_2(\eta)}{2d}}\,
\mathcal W_{K_\varepsilon}
\underset{\varepsilon\to0}{\longrightarrow}
\Wass_2.
\]
\paragraph{Discrete Wasserstein distances on Markov chains.}\label{sec-discrete-wasserstein-markov}
\(\mathsf J_{ij}\eqdef \pi_iK_{ij}=\pi_jK_{ji}\)
\[
\simplex_n\eqdef\left\{a\in\RR_+^n:\sum_i a_i=1\right\}.
\]
\(\rho_i(a)\eqdef \frac{a_i}{\pi_i}, \qquad a_i=\pi_i\rho_i(a).\)
The logarithmic mean \(\theta\) defined in~\eqref{eq-logarithmic-mean} is the mobility selected so that the entropy calculus later recovers exactly the Markov evolution; see Proposition~\ref{prop-discrete-markov-entropy-gradient}. The identity \(\theta(a,b)(\log a-\log b)=a-b\) converts entropy gradients into density differences along graph edges.

\(\bar\nabla\psi(i,j)\eqdef\psi_j-\psi_i.\)
\begin{equation}
\label{eq-discrete-markov-onsager}
(\mathcal K_\rho\psi)_i
\eqdef
\sum_j K_{ij}\theta(\rho_i,\rho_j)(\psi_i-\psi_j),
\qquad
(\mathcal L_a\psi)_i
\eqdef
\pi_i(\mathcal K_{\rho(a)}\psi)_i.
\end{equation}
\begin{equation}
\label{eq-discrete-markov-action}
\mathbb A_K(a,\psi)
\eqdef
\frac12\sum_{i,j}\mathsf J_{ij}\theta(\rho_i(a),\rho_j(a))
(\bar\nabla\psi(i,j))^2.
\end{equation}
\begin{equation}
\label{eq-discrete-markov-distance}
\mathcal W_K^2(a_0,a_1)
\eqdef
\inf_{a_t,\psi_t}
\int_0^1\mathbb A_K(a_t,\psi_t)\,\d t,
\qquad
\dot a_t+\mathcal L_{a_t}\psi_t=0,
\end{equation}
\subsection{Dynamic Unbalanced Wasserstein Distances}
\label{sec-unbalanced-ot}
\paragraph{Balance equation and tangent variables.}
\[
\partial_t\rho_t+\nabla\!\cdot m_t=s_t,
\qquad
\int_0^1\!\int
\left(\frac{\norm{m_t}^2}{\rho_t}+\kappa^2\frac{s_t^2}{\rho_t}\right)\d x\,\d t,
\]
\paragraph{Reaction--transport action.}
\begin{equation}
A_\kappa(a,w,g)
\eqdef
a\bigl(\norm{w}^2+\kappa^2 g^2\bigr).
\label{eq-wfr-velocity-action}
\end{equation}
\begin{equation}
J_\kappa(a,m,r)
\eqdef
\begin{cases}
\displaystyle\frac{\norm{m}^2+\kappa^2 r^2}{a},
& a>0,\\[6pt]
0, & a=0,\ m=0,\ r=0,\\
+\infty, & a=0,\ (m,r)\neq(0,0).
\end{cases}
\label{eq-wfr-momentum-perspective}
\end{equation}
\begin{equation}
\mathbb J_\kappa(\alpha,\omega,\sigma)
\eqdef
\int
J_\kappa\!\left(
\frac{\d\alpha}{\d\lambda},
\frac{\d \omega}{\d\lambda},
\frac{\d \sigma}{\d\lambda}
\right)\d\lambda.
\label{eq-wfr-measure-action}
\end{equation}
\paragraph{Static and dynamic viewpoints.}
The dynamic balance-equation formula is not a separate object from static unbalanced OT: it is the least-action representation of the same cone distance. To make the normalization explicit, define on the cone $\mathfrak C[\RR^d]$ the squared cost

\begin{equation}\label{eq-wfr-scaled-cone-metric}
\Delta_\kappa\big((x,r),(y,s)\big)^2
\eqdef
4\kappa^2
\left[
r^2+s^2-2rs
\cos\!\left(
\frac{\norm{x-y}}{2\kappa}\wedge\frac{\pi}{2}
\right)
\right].
\end{equation}
The radii encode masses through the weighted projection $\mathsf P_2$ defined in Section~\ref{sec-unbalanced}. Accordingly, set

\begin{equation}\label{eq-wfr-scaled-cone-value}
\CW_\kappa(\alpha_0,\alpha_1)
\eqdef
\inf_{\substack{\gamma\in\Mm_+(\mathfrak C[\RR^d]^2)\\
\mathsf P_2\gamma_1=\alpha_0,\ \mathsf P_2\gamma_2=\alpha_1}}
\int
\Delta_\kappa\big((x,r),(y,s)\big)^2
\d\gamma(x,r,y,s).
\end{equation}
For $\kappa=1/2$, this is exactly the normalization of the Hellinger--Kantorovich cone cost used in Section~\ref{sec-unbalanced}.

\begin{prop}[Static/dynamic equivalence for unbalanced OT]\label{prop-static-dynamic-unbalanced}
For nonnegative finite measures $\alpha_0,\alpha_1$ on $\RR^d$, the dynamic value
\begin{equation}\label{eq-dynamic-unbalanced-ot}
\WFR_\kappa^2(\alpha_0,\alpha_1)
\eqdef
\inf_{\substack{\partial_t\alpha_t+\nabla\cdot \omega_t=\sigma_t\\
\alpha_{t=0}=\alpha_0,\ \alpha_{t=1}=\alpha_1}}
\int_0^1
\mathbb J_\kappa(\alpha_t,\omega_t,\sigma_t)\,\d t
\end{equation}
equals the static cone value~\eqref{eq-wfr-scaled-cone-value}. Hence $\WFR_\kappa=\CW_\kappa^{1/2}$ is the geodesic distance generated by the balance-equation least-action problem.
\end{prop}
\begin{proof}
The cone construction turns variation of mass into radial motion and spatial transport into angular motion on $\mathfrak C[\RR^d]$. The normalization can be checked on a smooth cone path. If its base position is $x_t$, its cone radius is $r_t$, and the projected mass is $a_t=r_t^2$, then its relative growth rate is $g_t=\dot a_t/a_t=2\dot r_t/r_t$. The infinitesimal energy induced by~\eqref{eq-wfr-scaled-cone-metric} is therefore
\[
4\kappa^2\dot r_t^2+r_t^2\norm{\dot x_t}^2
=
a_t\left(\norm{\dot x_t}^2+\kappa^2g_t^2\right),
\]
which is precisely the velocity action~\eqref{eq-wfr-velocity-action}.

Applying the Benamou--Brenier theorem on the cone to the lifted endpoint measures now gives a dynamic least-action problem whose static value is $\CW_\kappa$. Projecting a cone curve back to the base space with the squared cone radius as weight produces a measure curve $\alpha_t$, a spatial flux measure $\omega_t$, and a source measure $\sigma_t$ satisfying the balance equation. The cone kinetic energy decomposes exactly into the three-variable perspective action $\mathbb J_\kappa$ in~\eqref{eq-dynamic-unbalanced-ot}. Conversely, any finite-action triple $(\alpha_t,\omega_t,\sigma_t)$ can be lifted, after relaxation, to a cone curve whose radial and spatial velocities realize the source and flux with the same action. This is the standard static/dynamic identification for the Hellinger--Kantorovich and Wasserstein--Fisher--Rao metrics.

The two infima are therefore equal; lower semicontinuity gives the general finite-measure statement from the smooth positive case.
\end{proof}
\paragraph{Balanced versus unbalanced interpolations.}
\subsection{Variational Mean Field Games}
\label{sec-variational-mean-field-games}
\paragraph{From individual control to a population equilibrium.}
\begin{equation}\label{eq-mfg-individual-control}
\inf_{\gamma(0)=x}
\left\{
\int_0^1
\left(
\frac12\norm{\dot\gamma(t)}^2
+g\bigl(\rho_t(\gamma(t))\bigr)
\right)\d t
+\Psi(\gamma(1))
\right\}.
\end{equation}
\paragraph{A Benamou--Brenier planning problem.}
\begin{equation}\label{eq-mfg-congestion-primitive}
g(r)=G'(r)
\end{equation}
\begin{equation}\label{eq-variational-mfg-velocity}
\inf_{(\rho_t,v_t)}
\left\{
\int_0^1\!\int_{\RR^d}
\left(
\frac12\rho_t(x)\norm{v_t(x)}^2+G(\rho_t(x))
\right)\d x\,\d t
+\int_{\RR^d}\Psi(x)\rho_1(x)\d x
\right\},
\end{equation}
\begin{equation}\label{eq-variational-mfg-continuity}
\partial_t\rho_t+\diverg(\rho_t v_t)=0,
\qquad
\rho_{t=0}\,\d x=\alpha_0.
\end{equation}
Unlike in~\eqref{eq:benamou-brenier}, the final measure is not prescribed. The terminal potential $\Psi$ selects desirable final states softly, while the integral of $G(\rho_t)$ penalizes crowded intermediate configurations.

\paragraph{Convex momentum formulation.}
As in the Benamou--Brenier problem, the density--velocity integrand and the constraint in~\eqref{eq-variational-mfg-velocity}--\eqref{eq-variational-mfg-continuity} are not jointly convex. Set $\omega_t=\alpha_t v_t$, or $m_t=\rho_t v_t$ when densities exist, and define the congestion functional

\begin{equation}\label{eq-mfg-congestion-functional}
\mathcal C_G(\alpha)
\eqdef
\begin{cases}
\displaystyle\int_{\RR^d}G(\rho(x))\d x,
& \alpha=\rho\,\d x,\\[4pt]
+\infty,&\text{otherwise}.
\end{cases}
\end{equation}
\begin{equation}\label{eq-variational-mfg-momentum}
\inf_{(\alpha_t,\omega_t)}
\left\{
\int_0^1
\left(
\frac12\mathbb J(\alpha_t,\omega_t)
+\mathcal C_G(\alpha_t)
\right)\d t
+\int_{\RR^d}\Psi\,\d\alpha_1
\right\},
\qquad
\partial_t\alpha_t+\diverg\omega_t=0,
\quad \alpha_{t=0}=\alpha_0.
\end{equation}
\begin{equation}\label{eq-variational-mfg-density-momentum}
\frac12J(\rho_t,m_t)+G(\rho_t)
=
\frac{\norm{m_t}^2}{2\rho_t}+G(\rho_t),
\end{equation}
\paragraph{Optimality system and game interpretation.}
\begin{prop}[Potential MFG system]\label{prop-variational-mfg-system}
Assume that an optimizer of~\eqref{eq-variational-mfg-momentum} has a smooth positive density $\rho$, that $G$ is differentiable, and set $g=G'$. Then there is a value function $u$ such that
\begin{equation}\label{eq-variational-mfg-system}
\left\{
\begin{aligned}
-\partial_tu+\frac12\norm{\nabla u}^2&=g(\rho),
& u_{t=1}&=\Psi,\\
\partial_t\rho-\diverg(\rho\nabla u)&=0,
& \rho_{t=0}\,\d x&=\alpha_0,
\end{aligned}
\right.
\qquad
m=-\rho\nabla u.
\end{equation}
\end{prop}
\begin{proof}
Use $-u$ as the multiplier for $\partial_t\rho+\diverg m=0$. After integration by parts, the space--time Lagrangian density is \(\frac{\norm m^2}{2\rho}+G(\rho) +\rho\,\partial_tu+m\cdot\nabla u.\)
Stationarity in $m$ gives $m=-\rho\nabla u$. Stationarity in $\rho$, followed by this substitution, gives the first equation in~\eqref{eq-variational-mfg-system}. The free terminal density contributes $(\Psi-u_1)\rho_1$, hence $u_1=\Psi$ under the positivity assumption. Primal feasibility gives the second equation.
\end{proof}
\paragraph{Hard congestion through a bottleneck.}
\begin{equation}\label{eq-variational-mfg-hard-cap}
G_\kappa(r)=\iota_{[0,\kappa]}(r)
=
\begin{cases}
0,&0\leq r\leq\kappa,\\
+\infty,&\text{otherwise}.
\end{cases}
\end{equation}
\begin{equation}\label{eq-variational-mfg-quadratic-terminal}
\Gamma_\eta(\rho_1)
=
\frac{\eta}{2}
\int_\Omega\abs{\rho_1(x)-\rho_\star(x)}^2\d x,
\end{equation}
where $\rho_\star$ is a desired terminal density.

The role of endpoint relaxation depends on the geometry. If two fixed endpoint densities on $\RR^d$ satisfy a spatially constant cap, their ordinary $\Wass_2$ geodesic satisfies the same cap by Proposition~\ref{prop-density-cap-geodesic-convex}, so the cap does not alter that geodesic.

\paragraph{Computational consequence.}